\section{Dendriform splitting and a non-special pre-Novikov algebra}

Another example of a variety
$\Var $ not satisfying the conditions
of Theorem~\ref{thm:Weight-Embedding}
is the class $\Zinb$ of Zinbiel (dual Leibniz or pre-commutative) algebras.
This is a particular case of the dendriform
splitting of a binary operad described in
\cite{BBGN, GubKol2014}.
Namely,
if $\Var $ is a variety of algebras with (one or more) binary operation
$\mu (x,y) = xy$ satisfying a family
of multi-linear identities $\Sigma $
then
$\pre\Var $ is a variety of algebras with
duplicated set of binary operations
$\mu_\vdash (x,y)=x\vdash y $,
$\mu_\dashv (x,y) = x\dashv y$
satisfying a set of identities $\pre\Sigma $
defined as follows.
Assume $f = f(x_1,\dots, x_n)$ is a multi-linear polynomial of degree $\le n$, and let $k\in \{1,\dots, n\}$.
Suppose $u$ is a nonassociative monomial in the variables
$x_1,\dots, x_n$ such that each $x_i$ appears in $u$
no more than once.
Define a polynomial~$u^{[k]}$ in~$x_1,\dots, x_n$
relative to the operations $\mu_\vdash $, $\mu _\dashv$
by induction on the degree.
If $u=x_i$ then
$u^{[k]} = x_i$;
if $u = vw$ and $x_k$ appears in $v$ (or in $w$) then
$u^{[k]} = v^{[k]}\dashv w^{[k]}$ (or, respectively,
$v^{[k]}\vdash w^{[k]}$);
if $x_k$ does not appear in $u$ then set
$u^{[k]}=
v^{[k]}\dashv w^{[k]} + v^{[k]}\vdash w^{[k]}$.
Transforming each monomial $u$ in the
multi-linear polynomial $f$
in this way, we get $f^{[k]}(x_1,\dots, x_n)$.
The collection of all such $f^{[k]}$
for $f\in \Sigma $, $k=1,\dots,\deg f$,
forms the set of defining relations of a new variety denoted
$\pre\Var $.

For example, for $f(x_1,x_2,x_3)
= (x_1x_2)x_3 - x_1(x_2x_3)$
the polynomials $f^{[k]}$, $k=1,2,3$, are given by
%
\begin{gather}
f^{[1]} = (x_1\dashv x_2)\dashv x_3 - x_1\dashv (x_2\dashv x_3 + x_2\vdash x_3),\nonumber \\
f^{[2]}
= (x_1\vdash x_2)\dashv x_3 - x_1\vdash (x_2\dashv x_3),\nonumber\\
f^{[3]}
= (x_1\vdash x_2 + x_1\dashv x_2)\vdash x_3 - x_1\vdash (x_2\vdash x_3).\label{eq:preAss}
\end{gather}
These identities define the variety of pre-associative
or dendriform algebras \cite{Loday}.

If the initial operation was commutative or anti-commutative then the set of identities $\pre\Sigma $
includes $x_1\vdash x_2 = \pm x_2\dashv x_1$,
so the operations in $\pre\Var $ are actually
expressed via $\mu_\vdash $ or $\mu_\dashv $.
For example, $\Var=\Lie$ produces the variety
$\pre\Lie$ of left- or right-symmetric algebras
(depending on the choice of $\vdash $ or $\dashv$).
If $\Var = \Com$ then, in terms of the operation
$x\cdot y = x\dashv y = y\vdash x$,
all three identities \eqref{eq:preAss}
of pre-associative algebras are equivalent to~\eqref{eq:Zinbiel}.

In a similar way, one may derive the identities
of a $\pre\Nov$-algebra by means of the dendriform
splitting applied to \eqref{eq:LSym} and \eqref{eq:RCom}.
Routine simplification leads us to the following
definition: a $\pre\Nov$-algebra is a linear space with
two bilinear operations $\vdash $, $\dashv$ satisfying
%
\begin{gather}
 (x_1\dashv x_2)\dashv x_3 = (x_1\dashv x_3)\dashv x_2, \nonumber\\
 (x_1\vdash x_2)\dashv x_3 = (x_1\vdash x_3)\vdash x_2 + (x_1\dashv x_3)\vdash x_2, \nonumber\\
 (x_1\dashv x_2)\dashv x_3 - x_1\dashv (x_2\dashv x_3) - x_1\dashv (x_2\vdash x_3)
 =
 (x_2\vdash x_1)\dashv x_3 - x_2\vdash (x_1\dashv x_3),
 \nonumber\\
 (x_1\vdash x_3)\dashv x_2 - x_1\vdash (x_2\vdash x_3)
 =
 (x_2\vdash x_3)\dashv x_1 - x_2\vdash (x_1\vdash x_3).\label{eq:pre-Novikov-2}
\end{gather}

The formal change of operations $x\dashv y = x\prec y$, $x\vdash y = y\succ x$
turns \eqref{eq:pre-Novikov-2} exactly into \eqref{eq:pre-Novikov}.
Hence, the operad $\pre\Nov = \pre D\Com$ defines
the same class of algebras as $D\Zinb = D\pre\Com $.

\begin{Remark}
This is not hard to compute that
$\pre D\Lie = D\pre\Lie$ and $\pre D\As = D\pre\As$.
In general, for every binary operad $\Var $ there exists
a morphism of operads $\pre D\Var \to D\pre\Var $
(i.e., every $D\pre\Var $-algebra is a $\pre D\Var$-algebra).
We do not know an example when
this morphism is not an isomorphism, i.e., when the
operations $\pre$ and $D$ applied to a binary operad
do not commute.
\end{Remark}

An equivalent way to define the variety
$\pre\Var $ was proposed in \cite{GubKol2014}.
Let $\Perm $ stand for the variety of associative
algebras that satisfy left commutativity
\[
x_1x_2x_3 = x_2x_1x_3.
\]
An algebra $V$ with two operations $\dashv$, $\vdash $
is a $\pre\Var$-algebra if and only if
for every $P\in \Perm $ the space
$P\otimes V$ equipped with the single operation
%
\begin{equation}\label{eq:Dend-Perm}
(p\otimes u)(q\otimes v) = pq\otimes (u\vdash v) +
qp\otimes (u\dashv v),
\qquad
p,q\in P,\quad u,v\in V,
\end{equation}
%
is a $\Var $-algebra.
The same statement holds in the case when the binary operad
$\Var $ is generated by several operations.

\begin{Remark}
For an arbitrary binary operad,
there is a morphism of operads
$\zeta\colon \pre\Var \to \Zinb\circ \Var $. Namely,
for every $A\in \Var $ and for every $Z\in \Zinb$
the space $Z\otimes A$ equipped with two operations
\[
(z\otimes a)\vdash (w\otimes b) = (w\cdot z)\otimes ab,
\qquad
(z\otimes a)\dashv (w\otimes b) = (z\cdot w)\otimes ab,
\]
for $z,w\in Z$, $a,b\in A$, is a $\pre\Var $-algebra.

However, $\pre\Var $ and $\Zinb\circ \Var $ are
not necessarily isomorphic. For example,
if $\Var $ is defined by the identity $(x_1\cdot x_2)\cdot x_3=0$
then the kernel of $\zeta $ is nonzero.
\end{Remark}

As a corollary, we obtain

\begin{Proposition}%\label{prop:preNov-Zinb}
The operad $\pre\Nov $
is isomorphic to the Manin white product $\Zinb\circ \Nov$.
\end{Proposition}

\begin{Remark}
Let $V$ be a $D\Zinb$-algebra.
In terms of pre-Novikov operations $\dashv$ and $\vdash$, the conformal algebra structure mentioned in
Remark~\ref{rem:Conformal} is expressed as
\[
(u_{(\lambda)} v) = \partial (v\dashv u) + \lambda (u\vdash v + v\dashv u),
\qquad u,v\in V.
\]
This is indeed a left-symmetric conformal algebra which is easy
to check via the conformal analogue of \eqref{eq:Dend-Perm}.
By slight abuse of notations,
for every $\Perm $-algebra $P$ the operation
%
\begin{gather*}
[(p\otimes u)_{(\lambda) } (q\otimes v)]
= pq \otimes (u_{(\lambda )} v) - qp\otimes (v_{(-\partial-\lambda )} u) \\
\qquad{}=pq\otimes (\partial (v\dashv u) + \lambda (u\vdash v + v\dashv u))
-qp\otimes (\partial (u\dashv v) - (\lambda+\partial) (v\vdash u + u\dashv v))\\
\qquad{}=\partial (qp\otimes v\vdash u+pq\otimes v\dashv u)
+\lambda (pq\otimes u\vdash v +qp\otimes u\dashv v
+ qp\otimes v\vdash u + pq\otimes v\dashv u)
\end{gather*}
%
is exactly the quadratic Lie conformal algebra structure \cite{Xu1999}
on $\Bbbk [\partial ]\otimes P\otimes V$
corresponding to the Novikov algebra $P\otimes V$:
\[
[x_{(\lambda )} y] = \partial (y x) + \lambda (x y + y x)
\]
for $x=p\otimes u$, $y=q\otimes v$, and the product
is given by \eqref{eq:Dend-Perm}.

Hence the construction of a $\pre\Lie$ conformal algebra from a
$D\Zinb $-algebra is a quite clear consequence of
the commutativity of tensor product.
\end{Remark}

The final statement of this section shows a substantial difference
between the properties of Novikov algebras and $\pre\Nov $-algebras.
Although the defining identities of $\pre\Nov $
are exactly those that hold on differential Zinbiel algebras
with operations \eqref{eq:Der-op}
(i.e.,
the dendriform analogue of \cite[Theorem~7.8]{DzhLofwall2002}
holds),
the general embedding statement (i.e., the dendriform analogue
of \cite[Theorem~3]{BCZ-2017})
turns to be wrong.

\begin{Theorem}%\label{thm:ZinbExmp}
If the characteristic of the base field $\Bbbk $ is not~$2$ or~$3$
then there exists a $D\Zinb $-algebra which cannot be embedded into a differential Zinbiel algebra.
\end{Theorem}

\begin{proof}
Consider the free Zinbiel algebra $F$ generated
by
\[
\{a,b\}^{(\omega )} = \bigl\{a,b,a',b', \dots, a^{(n)}, b^{(n)}, \dots \bigr\}.
\]
This is the free differential Zinbiel algebra
with two
generators $a$, $b$, its derivation $d$ maps $x^{(n)}$ to $x^{(n+1)}$
for $x\in \{a,b\}$. The product of two elements $f,g\in F$ is denoted
$f\cdot g$.

For every $f,g\in F$, define $f\prec g$, $f\succ g$
by the rule \eqref{eq:Der-op}:
\[
f\prec g = f\cdot d(g),\qquad f\succ g = d(f)\cdot g.
\]
Then $(F,\prec, \succ )$ is a $D\Zinb $-algebra, and its
subalgebra generated by $a$, $b$ is isomorphic to the free algebra $D\Zinb \langle a,b\rangle $.

Denote $f=b\prec b=b\cdot b' \in D\Zinb \langle a,b\rangle \subset F$, and let $J$ stand for the ideal in $F$ generated by~$f$ and all its derivatives:
\[
J = \bigl(f,f',f'', \dots \bigr) \triangleleft F,\qquad d(J)\subseteq J.
\]
In particular,
\[
h = a\cdot \bigl(f'\cdot b'\bigr)-a\cdot \bigl(f\cdot b''\bigr) \in J.
\]
Let us show that $h\in D\Zinb \langle a,b\rangle \subset J$.
Indeed,
\begin{align*}
h &= a\cdot \bigl(\bigl(b\cdot b'\bigr)'\cdot b'\bigr)-a\cdot \bigl(\bigl(b\cdot b'\bigr)\cdot b''\bigr) \\
 &= a\cdot \bigl(\bigl(b'\cdot b'\bigr)\cdot b'\bigr) + a\cdot \bigl(\bigl(b\cdot b''\bigr)\cdot b'\bigr)-a\cdot \bigl(\bigl(b\cdot b'\bigr)\cdot b''\bigr)
 = a\cdot \bigl(\bigl(b'\cdot b'\bigr)\cdot b'\bigr)
\end{align*}
due to the right commutativity of Zinbiel algebras.
Next,
$(b'\cdot b')\cdot b' = 2b'\cdot (b'\cdot b')$,
so
$(b'\cdot b')\cdot b'+b'\cdot (b'\cdot b')
= \frac{3}{2} (b'\cdot b')\cdot b'$.
Therefore,
\begin{align*}
h &= a\cdot \bigl(\bigl(b'\cdot b'\bigr)\cdot b'\bigr)
= \frac{2}{3} a\cdot \bigl(\bigl(b'\cdot b'\bigr)\cdot b' + b'\cdot \bigl(b'\cdot b'\bigr)\bigr)
 = \frac{2}{3} \bigl(a\cdot \bigl(b'\cdot b'\bigr)\bigr)\cdot b'
 \\
 &= \frac{1}{3} \bigl(\bigl(a\cdot b'\bigr)\cdot b'\bigr)\cdot b'= \frac{1}{3} ((a\prec b)\prec b)\prec b = 2\bigl[ab'b'b'\bigr]
 \in D\Zinb \langle a,b\rangle .
\end{align*}
As in Example~\ref{exmp:Zinbiel}, we denote by
$[x_1x_2\cdots x_{n-1}x_n]$ the following expression
in a Zinbiel algebra:
\[
[x_1x_2\cdots x_{n-1}x_n]
= x_1\cdot (x_2\cdot (\cdots (x_{n-1}\cdot x_n)\cdots )).
\]
Recall \cite{Loday} that all such expressions with $x_i$
from a set $X$ form a linear basis of the free Zinbiel algebra
generated by $X$ (i.e., this is a normal form in $\pre\Com\langle X\rangle $).

Let $I$ be the ideal in $D\Zinb \langle a,b\rangle $
generated by $f$, and let $V = D\Zinb \langle a,b\mid f\rangle
=D\Zinb \langle a,b\rangle/I $.
Then
$F/J$ is the universal differential Zinbiel envelope
of $V$. To prove the theorem, it remains to show that $h\notin I$:
if so then $h+I$
would lie in the kernel of every homomorphism
from $V$ to the derived algebra $Z^{(d)}$
constructed
from a differential Zinbiel algebra $Z$ with a derivation~$d$.

Assume $[ab'b'b']\in I$.
The specific of the Zinbiel identity \eqref{eq:Zinbiel}
is that the first letter remains unchanged in all terms.
Hence, $[ab'b'b']$ should be a
linear combination of the elements
$(a\ast f\star b)$, $(a\ast b\star f)$,
where $\ast, \star \in \{\prec, \succ \}$,
with two possible bracketing each, so we have in total 16 terms
under consideration. Let us write them all in the normal form in $F$:
\begin{gather*}
(a\prec f)\prec b
%= (a(bb')')b' = (ab')(bb')' = (ab')(b'b' + bb'')
 = 3\bigl[ab'b'b'\bigr]+\bigl[ab'bb''\bigr]+\bigl[abb'b''\bigr]+\bigl[abb''b'\bigr], \\
%
(a\prec f)\succ b
 % = (a(bb')')'b = (a'(bb')')b+(a(bb')'')b
 % = (a'b)(b'b'+bb'') + (ab)(b''b'+2b'b''+bb''')
 =\bigl[a'bb'b'\bigr]+\bigl[a'b'bb'\bigr]+\bigl[a'b'b'b\bigr]+2\bigl[a'bbb''\bigr]+\bigl[a'bb''b\bigr]+\bigl[abb''b'\bigr]+\bigl[ab''bb'\bigr]\\
\hphantom{(a\prec f)\succ b=}{} +\bigl[ab''b'b\bigr]+2\bigl[abb'b''\bigr]+2\bigl[ab'bb''\bigr]+2\bigl[ab'b''b\bigr]+2\bigl[abbb'''\bigr]+\bigl[abb'''b\bigr], \\
%
(a\succ f)\prec b
 %= (a'(bb'))b' = (a'b')(bb')
 = \bigl[a'b'bb'\bigr]+2\bigl[a'bb'b'\bigr], \\
%
(a\succ f)\succ b
 %= (a'(bb'))'b = (a''b)(bb') + (a'b)(b'b'+bb'')
 = 2\bigl[a''bbb'\bigr]+\bigl[a''bb'b\bigr] +\bigl[a'bb'b'\bigr]+\bigl[a'b'bb'\bigr]+\bigl[a'b'b'b\bigr]+2\bigl[a'bbb''\bigr]+\bigl[a'bb''b\bigr], \\
%
a\prec (f\prec b)
 %= a((bb')b')' = 2a\bigl[bb'b'\bigr]'
 = 2\bigl[ab'b'b'\bigr]+2\bigl[abb''b'\bigr]+2\bigl[abb'b''\bigr], \\
%
a\prec (f\succ b)
 %= a((bb')'b)' = a((b'b')b)'+a((bb'')b)'
 = a\bigl[b'b'b\bigr]'+a\bigl[b'bb'\bigr]' + a\bigl[bb''b\bigr]' + a\bigl[bbb''\bigr]'
 = \bigl[ab''b'b\bigr] + 2\bigl[ab'b''b\bigr]+ 2\bigl[ab'b'b'\bigr]\\
\hphantom{a\prec (f\succ b)=}{} + \bigl[ab''bb'\bigr]+2\bigl[ab'bb''\bigr] +\bigl[abb'''b\bigr]+\bigl[abb''b'\bigr]+\bigl[abb'b''\bigr]+\bigl[abbb'''\bigr], \\
%
a\succ (f\prec b)
 %= a'((bb')b')
 = 2\bigl[a'bb'b'\bigr], \\
%
a\succ (f\succ b)
 % = a'((bb')'b) = a'((b'b')b)+a'((bb'')b)
 = \bigl[a'b'b'b\bigr]+\bigl[a'b'bb'\bigr]+\bigl[a'bb''b\bigr]+\bigl[a'bbb''\bigr], \\
%
(a\prec b)\prec f
 % = (ab')(bb')' = (ab')(b'b')+(ab')(bb'')
 = 3\bigl[ab'b'b'\bigr] + \bigl[ab'bb''\bigr]+\bigl[abb'b''\bigr]+\bigl[abb''b'\bigr], \\
%
(a\prec b)\succ f
 % = (ab')'(bb') = (a'b')(bb')+(ab'')(bb')
 = \bigl[a'b'bb'\bigr]+2\bigl[a'bb'b'\bigr] + \bigl[ab''bb'\bigr]+\bigl[abb''b'\bigr]+\bigl[abb'b''\bigr], \\
%
(a\succ b)\prec f
 %= (a'b)(bb')' = (a'b)(b'b')+(a'b)(bb'')
 = \bigl[a'bb'b'\bigr]+\bigl[a'b'bb'\bigr]+\bigl[a'b'b'b\bigr] + 2\bigl[a'bbb''\bigr]+\bigl[a'bb''b\bigr], \\
%
(a\succ b)\succ f
 %= (a'b)'(bb') = (a''b)(bb') + (a'b')(bb')
 = 2\bigl[a''bbb'\bigr]+\bigl[a''bb'b\bigr]+\bigl[a'b'bb'\bigr]+2\bigl[a'bb'b'\bigr], \\
%
a\prec (b\prec f)
 % = a(b(bb')')' = a\bigl[bb'b'\bigr]' + a\bigl[bbb''\bigr]'
 = \bigl[ab'b'b'\bigr]+\bigl[abb''b'\bigr] +2\bigl[abb'b''\bigr] +\bigl[ab'bb''\bigr] + \bigl[abbb'''\bigr], \\
%
a\prec (b\succ f)
 %= a(b'(bb'))' = a\bigl[b'bb'\bigr]'
 = \bigl[ab''bb'\bigr]+\bigl[ab'b'b'\bigr]+\bigl[ab'bb''\bigr], \\
%
a\succ (b\prec f)
 %= a'(b(bb')') = a'\bigl[bb'b'\bigr]+a'\bigl[bbb''\bigr]
 = \bigl[a'bb'b'\bigr] + \bigl[a'bbb''\bigr], \\
%
a\succ (b\succ f)
 %= a'(b'(bb'))
 = \bigl[a'b'bb'\bigr].
\end{gather*}
Arrange the normal Zinbiel words in the following order:
$[ab'b'b']$, $[ab'bb'']$, $[abb'b'']$, $[abb''b']$, $[ab''bb']$, $[ab''b'b]$, $[ab'b''b]$, $[abbb''']$, $[abb'''b]$, $[a''bb'b]$, $[a''bbb']$, $[a'bb'b']$, $[a'b'bb']$, $[a'b'b'b]$, $[a'bbb'']$, $[a'bb''b]$.
Then the assumption $[ab'b'b']\in I$ is equivalent to the condition
that the row vector $e_1 = (1,0,\dots ,0)\in \Bbbk ^{16}$
belongs to the row space of the matrix
%
\[ \left(
\begin{array}{cccccccccccccccc}
3& 1& 1& 1& 0& 0& 0& 0& 0& 0& 0& 0& 0& 0& 0& 0\\
0& 2& 2& 1& 1& 1& 2& 2& 1& 0& 0& 1& 1& 1& 2& 1\\
0& 0& 0& 0& 0& 0& 0& 0& 0& 0& 0& 2& 1& 0& 0& 0\\
0& 0& 0& 0& 0& 0& 0& 0& 0& 1& 2& 1& 1& 1& 2& 1\\
2& 0& 2& 2& 0& 0& 0& 0& 0& 0& 0& 0& 0& 0& 0& 0\\
2& 2& 1& 1& 1& 1& 2& 1& 1& 0& 0& 0& 0& 0& 0& 0\\
0& 0& 0& 0& 0& 0& 0& 0& 0& 0& 0& 2& 0& 0& 0& 0\\
0& 0& 0& 0& 0& 0& 0& 0& 0& 0& 0& 0& 1& 1& 1& 1\\
3& 1& 1& 1& 0& 0& 0& 0& 0& 0& 0& 0& 0& 0& 0& 0\\
0& 0& 1& 1& 1& 0& 0& 0& 0& 0& 0& 2& 1& 0& 0& 0\\
0& 0& 0& 0& 0& 0& 0& 0& 0& 0& 0& 1& 1& 1& 2& 1\\
0& 0& 0& 0& 0& 0& 0& 0& 0& 1& 2& 2& 1& 0& 0& 0\\
1& 1& 2& 1& 0& 0& 0& 1& 0& 0& 0& 0& 0& 0& 0& 0\\
1& 1& 0& 0& 1& 0& 0& 0& 0& 0& 0& 0& 0& 0& 0& 0\\
0& 0& 0& 0& 0& 0& 0& 0& 0& 0& 0& 1& 0& 0& 1& 0\\
0& 0& 0& 0& 0& 0& 0& 0& 0& 0& 0& 0& 1& 0& 0& 0
\end{array} \right).
\]
This matrix may be transformed by elementary row and column
transformations with integer coefficients
to the triangular form with $1$, $2$, or $0$ on the diagonal,
so that its rank equals 10 for $\operatorname{char}\Bbbk \ne 2$.
Adding one more row $e_1$ increases the rank, so $[ab'b'b']\notin I$.
\end{proof}
